\documentclass[11pt]{article}

%— Margins & encoding
\usepackage[T1]{fontenc}
\usepackage[utf8]{inputenc}
\usepackage[margin=1in]{geometry}
\usepackage{setspace}

%— Math, graphics, tables, citations
\usepackage{amsmath,amssymb}
\usepackage{graphicx}
\usepackage{booktabs}
\usepackage{siunitx}
\usepackage{natbib}
\usepackage[hidelinks]{hyperref}
\usepackage{enumitem}
\usepackage{microtype}
\usepackage{xcolor}

%— Spacing
\onehalfspacing

\DeclareUnicodeCharacter{2006}{\,}
\DeclareUnicodeCharacter{202F}{\,}

%— Title block
\title{\large\textbf{Momentum in Large-Cap Equities:\\
Does the Classic 12--1 Strategy Still Work?}\\[0.5em]
\normalsize A Survivorship-Bias-Free Examination of the S\&P~500, 2006--2024}

\author{%
  Darshan Sathish Kumar\thanks{%
    Email: \href{mailto:darshansathishkumar@gmail.com}{darshansathishkumar@gmail.com}.
    Code and replication data: \url{https://github.com/dshan12/Momentum-Research}.
  }
}
\date{\today}

% ------------------------------------------------------------------ %
\begin{document}
\maketitle

% ------------------------------------------------------------------ %
\begin{abstract}
\noindent
We evaluate the canonical 12--1 month cross-sectional momentum strategy applied to the
S\&P~500 over the period March~2006--December~2024, using a survivorship-bias-free
constituent history, gap-aware return construction, and a granular turnover-based
transaction cost model. Each month we rank active constituents by their prior 12-month
return (skipping the most recent month), go long the top decile and short the bottom
decile equal-weighted, and rebalance monthly. The strategy produces a net annualized
return of $-2.79\%$, annualized volatility of $20.66\%$, a Sharpe ratio of $-0.23$,
and a maximum drawdown of $-81\%$ at 10~basis points per side. Leg decomposition
reveals that the signal is economically valid---the long leg earns $+7.9\%$ annualized
and the short leg loses $-9.1\%$---but systematic shorting of past losers proves
catastrophically unprofitable during the 2009 and 2020 momentum crash episodes, when
deeply distressed stocks recovered by $40$--$100\%$ within single calendar months.
Average monthly one-way turnover is $53.6\%$; transaction costs account for only
$\approx0.05\%$ per month and are not the primary driver of losses. Factor regressions
(CAPM through FF5$+$UMD, with Newey--West errors) show that $R^2$ rises from $0.15$
under CAPM to $0.82$ once UMD is included, with a statistically significant negative
alpha of $-4.0\%$ annually ($t = -3.23$). These results confirm that na\"ive large-cap
momentum is primarily a leveraged exposure to the well-documented UMD factor, with no
residual alpha once crash risk and frictions are accounted for, and that the strategy's
poor performance over this period is driven by the short leg rather than by any failure
of the ranking signal itself.
\end{abstract}

\noindent\textbf{Keywords:} momentum; cross-sectional strategies; momentum crashes;
survivorship bias; transaction costs; Fama--French factors; S\&P~500.\\
\textbf{JEL codes:} G11; G12; C58.

% ================================================================== %
\section{Introduction}
% ================================================================== %

The momentum anomaly---the empirical regularity that stocks with strong recent returns
tend to continue outperforming those with weak recent returns---has occupied a central
place in the asset pricing literature since \citet{Jegadeesh1993} first documented it
systematically in U.S.\ equities. Three decades on, momentum remains simultaneously one
of the most replicated and most debated findings in finance. It appears across asset
classes \citep{Asness2013}, geographies \citep{Rouwenhorst1998}, and time periods
stretching back to Victorian-era stock markets \citep{Geczy2013}. Yet it is also a
strategy that has spectacularly failed practitioners at precisely the worst moments
\citep{Daniel2016}, and whose theoretical foundations remain contested
\citep{Barberis1998,Hong1999}.

The U.S.\ large-cap equity universe---specifically the S\&P~500---presents a
particularly challenging testing ground. Constituent stocks are among the most
liquid, most widely followed, and most efficiently priced securities in the world.
Any systematic pattern in their returns attracts immediate arbitrage capital, compressing
risk premia rapidly. Moreover, the post-2000 era has been unusually hostile to momentum
strategies: two of the most severe momentum crashes on record---March--May 2009 and
November~2020---occurred within this window, and the prolonged bull market of 2010--2024
created conditions under which past losers persistently recovered, undermining the
short leg of momentum portfolios.

\paragraph{Research question.}
\emph{Can a na\"ive 12--1 month momentum rule, applied to S\&P~500 constituents with
survivorship-bias-free data and realistic transaction costs, generate statistically
significant abnormal returns over the 2006--2024 period? And if not, which component
of the strategy---the long leg, the short leg, or transaction costs---is responsible
for its failure?}

\paragraph{Contribution.}
This paper makes four contributions relative to prior replication work. First, we
reconstruct a dated S\&P~500 membership history from Wikipedia change logs, enabling
gap-aware return computation that avoids both survivorship bias (including only
current constituents) and re-entry bias (computing returns only across consecutive
months in which a stock is an active constituent). Second, we implement a granular
turnover-based cost model that tracks mark-to-market weight drift between rebalancing
dates, producing realistic one-way turnover estimates of $\approx54\%$ per month.
Third, we decompose gross P\&L into long-leg and short-leg contributions, providing
a clearer diagnosis of where the strategy fails than aggregate performance statistics
alone. Fourth, we contextualize our results against the Fama--French UMD factor
over the same sample period, demonstrating that the strategy's return profile closely
tracks UMD---including its crash behavior---while underperforming it net of costs.

\paragraph{Preview of results.}
The signal works as a cross-sectional ranking mechanism: the long leg earns $+7.9\%$
annualized and the short leg loses $-9.1\%$ in absolute terms, confirming that past
winners outperform past losers on average. However, the combined strategy is unprofitable
because the short leg bears catastrophic losses during two momentum crash episodes that
together erase more than half the portfolio's value. Transaction costs, while non-trivial
at $\approx0.65\%$ annualized, are a secondary concern. The strategy's return variation
is almost entirely spanned by the UMD factor ($R^2 = 0.82$), and its alpha is negative
and significant once UMD is controlled for. These findings reinforce the view that
momentum in large-cap U.S.\ equities is a priced risk factor---specifically, crash
risk---rather than a persistent source of exploitable alpha.

% ================================================================== %
\section{Related Literature}
\label{sec:lit}
% ================================================================== %

\citet{Jegadeesh1993} establish the foundational result: strategies that buy past
six- to twelve-month winners and sell past losers earn abnormal returns of approximately
$1\%$ per month over a three- to twelve-month holding period in U.S.\ equities.
\citet{carhart1997persistence} show that momentum largely explains the persistence of
mutual fund performance, cementing it as a systematic risk factor. \citet{Fama1996}
incorporate momentum into their three-factor framework's ``biggest embarrassment,''
acknowledging that their model fails to capture it; the formal UMD factor is later
constructed to address this gap \citep{FamaFrench2015}.

The theoretical literature offers two broad classes of explanations. Behavioral models
attribute momentum to investor underreaction \citep{Barberis1998} or gradual information
diffusion \citep{Hong1999}: prices adjust too slowly to fundamental news, creating
predictable return continuation. Risk-based explanations posit that momentum portfolios
earn compensation for bearing some form of systematic risk, though identifying the precise
risk factor has proven elusive \citep{Liu2002,Johnson2002}.

A second strand of literature documents the limits and costs of momentum strategies.
\citet{NovyMarx2016} demonstrate that momentum strategies require substantially higher
turnover than previously appreciated and that realistic transaction costs---particularly
for institutional investors trading at scale---can eliminate much of the gross alpha.
\citet{Frazzini2015} estimate effective trading costs for a large institutional manager
and find that momentum remains profitable after costs only when carefully implemented
with patient execution. \citet{Lesmond2004} show that much of the gross momentum premium
accrues in illiquid stocks where bid-ask spreads are prohibitive, calling into question
the economic significance of the anomaly for large portfolios.

Perhaps most relevant to our findings is the literature on momentum crashes.
\citet{Daniel2016} demonstrate that momentum strategies have highly negatively skewed
return distributions, with occasional catastrophic losses (``crashes'') that occur
during market rebounds following sharp downturns. The mechanism is straightforward:
momentum strategies accumulate large short positions in stocks that have fallen
dramatically; when the market reverses, these stocks exhibit the highest betas and
recover most violently, inflicting concentrated losses on the short book. The authors
document monthly losses of $-49\%$ (March--May 1932), $-49\%$ (March--May 2009), and
other severe episodes. \citet{Barroso2015} show that volatility-scaling the momentum
position---reducing exposure when recent strategy volatility is elevated---substantially
improves risk-adjusted performance and Sharpe ratios by approximately 50\%.

More recent work examines the performance of momentum in the post-financial-crisis
period. \citet{Asness2019} document that momentum has delivered positive returns over
very long samples (1926--present) but with significant time-variation and a notable
deterioration post-2000 in large-cap equity markets. \citet{Gulen2021} argue that
increased arbitrage activity and factor crowding have compressed the momentum premium
in large-cap stocks specifically, while it remains robust in smaller-cap universes.

Our paper contributes to this literature by providing a clean, reproducible implementation
with full code transparency, a focus on the S\&P~500 universe over a recent and economically
challenging period, and an explicit decomposition that isolates the short leg as the source
of losses rather than transaction costs or signal failure.

% ================================================================== %
\section{Data}
\label{sec:data}
% ================================================================== %

\subsection{Universe and Membership History}

The investment universe consists of all firms that appeared as S\&P~500 constituents
at any point between January~2005 and December~2024. We identify $826$ unique tickers
over this window. Historical additions and deletions are reconstructed from the Wikipedia
S\&P~500 constituent change log, which records the date, added ticker, and removed ticker
for each index rebalancing event since the 1990s. Starting from the current index
composition, we roll backward through all change events to reconstruct the constituent
set at each calendar month-end, then forward to produce a monthly membership panel
spanning the full sample. Each ticker is marked as active only in months in which it
was an S\&P~500 constituent, ensuring that returns are computed only for stocks that
were genuinely investable at the relevant time.

This construction eliminates survivorship bias in both directions: stocks deleted from
the index (due to merger, bankruptcy, or downgrade) contribute their returns up to
deletion and are then excluded, preventing the artificial inflation of historical
performance from holding only survivors. The panel averages $393$ active constituents
per month (range: $0$--$502$), with lower coverage in the early sample period as the
Wikipedia history becomes less complete prior to 2005.

\subsection{Prices and Return Construction}

End-of-month adjusted closing prices are downloaded via \texttt{yfinance}
(\texttt{interval="1mo"}, \texttt{auto\_adjust=True}). All prices are snapped to
month-end dates to ensure alignment with the membership panel.

Monthly returns are constructed as:
\[
  r_{i,t} = \frac{p_{i,t}}{p_{i,t-1}} - 1,
\]
but only when \emph{both} $p_{i,t}$ and $p_{i,t-1}$ are non-missing and the stock is
an active constituent in month $t$. Critically, we do \emph{not} use \texttt{pct\_change}
with gap-filling, which would compute the return from the last available price across
an arbitrarily long gap. When a stock re-enters the index after an absence, its return
in the re-entry month is set to missing, preventing the accumulation of a multi-month
price move from being attributed to a single period. This is an important correction:
without it, stocks that crashed and recovered during an index absence can appear to have
single-month returns of $+100\%$ or more, severely distorting cross-sectional rankings
and portfolio returns.

A hard cap of $\pm40\%$ is applied to monthly individual stock returns prior to
winsorization. S\&P~500 constituents essentially never move more than $40\%$ in a
single calendar month under normal market conditions; values beyond this threshold
almost always reflect data artifacts (split adjustments, corporate actions) or
represent returns on deeply distressed stocks that would not be replicable in a leveraged
short book due to margin requirements and forced stop-outs.\footnote{Without this cap,
multiple financial and cyclical stocks in our short book recorded monthly returns of
$+97\%$--$+100\%$ in April~2009, producing a single-month strategy loss of $-53\%$.
With the $40\%$ cap, the April~2009 loss is $-31\%$, which remains the worst month in
the sample but is more consistent with the $-49\%$ cumulative loss documented for the
UMD factor over March--May 2009 \citep{Daniel2016}.} Following the hard cap, returns
are winsorized cross-sectionally at the 1st and 99th percentiles in months with at
least 50 active names.

\subsection{Factors and Benchmarks}

Monthly Fama--French five-factor returns (Mkt-RF, SMB, HML, RMW, CMA), the risk-free
rate (RF), and the momentum factor (UMD) are obtained from Kenneth French's data
library.\footnote{\url{https://mba.tuck.dartmouth.edu/pages/faculty/ken.french/data_library.html}}
All factor returns are expressed as monthly decimals. The U.S.\ market benchmark is
constructed as $\text{Mkt-RF} + R_f$, yielding $+12.0\%$ annualized over the
2006--2024 sample. The UMD factor produces $+0.7\%$ annualized over the same period,
consistent with the documented deterioration of the momentum premium in large-cap
U.S.\ equities post-2000.

\subsection{Descriptive Statistics}

Table~\ref{tab:summary} reports summary statistics for the net strategy returns and
the U.S.\ market over the 226-month sample period. The strategy exhibits substantially
higher volatility than the market ($20.66\%$ vs.\ $15.20\%$ annualized), a large
negative skew ($-1.21$ vs.\ $-0.55$), and excess kurtosis ($5.84$ vs.\ $1.40$),
reflecting the fat left tail generated by momentum crash months. The maximum drawdown
of $-81.2\%$ dwarfs that of the market ($-50.3\%$), driven by the 2009 crash from
which the strategy never fully recovered before encountering the 2020 crash.

\begin{table}[!ht]
\centering
\caption{Summary statistics. Monthly decimal returns; sample March~2006--December~2024
         ($n = 226$ months). Sharpe computed with $r_f = 2\%$ annualized.
         Ann Vol $= \sigma_m \times \sqrt{12}$.}
\label{tab:summary}
\begin{tabular}{lrrrrrrrr}
\toprule
Series & $n$ & Mean\,(m) & Sharpe & Ann\,Ret & Ann\,Vol & Skew & Kurt & Max\,DD \\
\midrule
Strategy (net) & 226 & $-$0.0026 & $-$0.23 & $-$2.79\% & 20.66\% & $-$1.21 & 5.84 & $-$81.2\% \\
US Market      & 226 & \phantom{$-$}0.0094 & \phantom{$-$}0.61 & \phantom{$-$}12.02\% & 15.20\% & $-$0.55 & 1.40 & $-$50.3\% \\
\bottomrule
\end{tabular}
\end{table}

% ================================================================== %
\section{Methodology}
\label{sec:method}
% ================================================================== %

\subsection{Signal Construction}

For each month $t$, the momentum signal for stock $i$ is the cumulative log-price
appreciation from $t{-}13$ to $t{-}1$:
\[
  \text{mom}_{i,t} \;=\; \ln P_{i,t-1} \;-\; \ln P_{i,t-13}.
\]
This is the standard 12--1 specification: a 12-month lookback window that skips the
most recent month. Skipping the most recent month is standard practice to avoid
contamination by short-term price reversals \citep{Jegadeesh1990}, which would partially
offset momentum at the one-month horizon. The signal is computed only for stocks with
valid (non-missing) prices in both $t{-}1$ and $t{-}13$; stocks with gaps in either
endpoint are excluded from the ranking in that month.

\subsection{Portfolio Formation}

At each month-end $t$, active constituents with valid momentum signals are ranked into
percentiles. The long portfolio holds stocks in the top decile ($\geq 90$th percentile)
and the short portfolio holds stocks in the bottom decile ($\leq 10$th percentile).
Both legs are equal-weighted within the available names. We require a minimum of 20
stocks on each leg before executing; months failing this threshold are treated as
no-trade months and excluded from the return series ($25$ of $252$ calendar months
are excluded on this basis). Weights formed at the end of month $t$ are applied to
returns realized during month $t{+}1$, ensuring no look-ahead bias.

\subsection{Transaction Cost Model}

Let $w_{i,t}$ denote the target weight for stock $i$ entering month $t$. The
\emph{pre-trade} weight at time $t$---what the portfolio holds just before
rebalancing---is the prior-period weight marked to market:
\[
  \tilde{w}_{i,t} \;=\; \frac{w_{i,t-1}(1 + r_{i,t-1})}
    {\displaystyle\sum_j \bigl|w_{j,t-1}(1 + r_{j,t-1})\bigr|},
\]
where the denominator renormalizes each leg independently to maintain dollar
neutrality.\footnote{Each leg is renormalized separately: the long leg sums to $+1$
and the short leg sums to $-1$, preserving the gross notional of 2.}
One-way turnover is then:
\[
  \text{TO}_t \;=\; \tfrac{1}{2} \sum_i \bigl| w_{i,t} - \tilde{w}_{i,t} \bigr|,
\]
and the net portfolio return is:
\[
  r^{\text{net}}_t \;=\; \sum_i w_{i,t-1}\,r_{i,t} \;-\; c \cdot \text{TO}_t,
\]
with one-way cost $c = 10$~bps. Average monthly one-way turnover is $53.6\%$ (median
$57.2\%$; 95th percentile $83.8\%$), implying average monthly costs of $\approx0.054\%$
at the baseline assumption. We evaluate sensitivity to cost assumptions from 5 to 25~bps
in Section~\ref{sec:results}.

\subsection{Risk-Adjusted Performance}

Abnormal performance is estimated by regressing monthly excess strategy returns on
standard factor models:
\[
  r^{\text{net}}_t - r_{f,t} \;=\; \alpha + \boldsymbol{\beta}^\top \mathbf{f}_t + \varepsilon_t,
\]
where $\mathbf{f}_t$ contains the relevant factor returns in a given specification
(CAPM: $\{\text{Mkt-RF}\}$; FF3: $+$SMB, HML; FF5: $+$RMW, CMA; FF5$+$UMD: $+$UMD).
All regressions use OLS with Newey--West heteroskedasticity-and-autocorrelation-consistent
standard errors with 6 lags, appropriate for the overlapping information structure of
monthly momentum signals.

% ================================================================== %
\section{Results}
\label{sec:results}
% ================================================================== %

\subsection{Aggregate Performance}

The strategy produces a gross annualized return of approximately $-2.4\%$ and a net
annualized return of $-2.79\%$ at 10~bps per side, indicating that transaction costs
($\approx0.65\%$ annualized) are a secondary drag rather than the primary source of
underperformance. Annualized volatility of $20.66\%$ substantially exceeds that of the
U.S.\ market ($15.20\%$), reflecting the strategy's net short equity exposure during
crash months when short-side losses dominate. The Sharpe ratio is $-0.23$ and the
maximum drawdown is $-81.2\%$, compared to $+0.61$ and $-50.3\%$ for the market.
Figure~\ref{fig:equity} displays cumulative wealth for the strategy, the U.S.\ market,
and the UMD factor from March~2006.

\begin{figure}[!ht]
  \centering
  \includegraphics[width=\textwidth]{figures/equity_curve.png}
  \caption{Cumulative wealth indexed to $1.0$ at March~2006.
           \emph{Strategy (net)}: long-short portfolio after 10~bps per-side costs.
           \emph{Strategy (gross)}: before costs.
           \emph{US Market}: $\text{Mkt-RF} + R_f$ from Ken French data library.
           \emph{UMD}: Fama--French momentum factor.
           The strategy peaks at $\approx1.6\times$ in mid-2008, collapses during the
           2009 momentum crash, and declines through the subsequent bull market as the
           short leg persistently loses money.}
  \label{fig:equity}
\end{figure}

The strategy's behavior is dominated by two extreme events. In April~2009, the strategy
lost $-31\%$ in a single month (capped from a raw $-53\%$): deeply distressed financial
and cyclical stocks---Wynn Resorts, Principal Financial, Office Depot, Ciena, and
others---which had been correctly identified as bottom-decile momentum stocks after
their 2008 collapse, recovered by $40$--$100\%$ as the market bottomed. In
November~2020, the strategy lost $-19\%$ as COVID-era losers surged on vaccine news.
These two months alone account for approximately $50\%$ of the total drawdown. The
drawdown trajectory (Figure~\ref{fig:dd}) makes clear that the strategy never recovered
from the 2009 crash before encountering 2020.

\begin{figure}[!ht]
  \centering
  \includegraphics[width=\textwidth]{figures/drawdown.png}
  \caption{Drawdown series for the net strategy (red) and the U.S.\ market (blue dashed),
           computed as the percentage decline from each series' running maximum.
           The strategy's $-81\%$ maximum drawdown reflects sequential crashes in 2009
           and 2020 from which there was no meaningful recovery in between.}
  \label{fig:dd}
\end{figure}

\subsection{Leg Decomposition}

A key diagnostic---and the central interpretive finding of this paper---is the
decomposition of gross performance into the contributions of the long and short legs.
Table~\ref{tab:legs} reports this decomposition.

\begin{table}[!ht]
\centering
\caption{Long-leg and short-leg performance decomposition. Gross monthly returns,
         sample March~2006--December~2024 ($n = 226$ months).}
\label{tab:legs}
\begin{tabular}{lrr}
\toprule
Leg & Mean monthly return & Annualized return \\
\midrule
Long (top decile, past winners)  & $+$0.634\% & $+$7.9\% \\
Short (bottom decile, past losers) & $-$0.792\% & $-$9.1\% \\
\midrule
Combined (long $-$ short) & $-$0.158\% & $-$1.9\% \\
\bottomrule
\end{tabular}
\end{table}

\noindent
These results contain a striking and important finding. The long leg earns $+7.9\%$
annualized, which is economically meaningful and consistent with the signal correctly
identifying stocks with above-average forward returns. The signal \emph{works} as a
ranking mechanism. The short leg, however, loses $-9.1\%$ annualized in absolute
terms---meaning that the stocks identified as bottom-decile momentum names go \emph{up}
on average, costing the portfolio money on the short side.

This is not a failure of the ranking signal: the short leg contains stocks with the
weakest prior returns, and they do underperform the long leg (the $+7.9\%$ vs.\ $-9.1\%$
spread of $17\%$ is the gross momentum premium). The problem is that past losers in the
S\&P~500 also tend to go up in absolute terms in this sample, particularly during the
two crash recovery episodes. The strategy is long the market on the long leg and short
the market on the short leg, but the short leg's market beta is far higher during
recovery months than during normal months---exactly the dynamic documented theoretically
by \citet{Daniel2016} as the source of momentum crashes.

The implication is important for practitioners: \emph{the problem with this strategy is
not that momentum does not work, but that the short leg is not investable in its naive
form within a large-cap, leveraged framework.} A long-only implementation---holding top-decile
momentum stocks as an active overweight against an index benchmark---would have earned
$+7.9\%$ gross over this period, likely with a positive Sharpe ratio and manageable drawdowns.

\subsection{Transaction Cost Sensitivity}

Table~\ref{tab:tc} reports performance across a range of one-way cost assumptions.
The strategy is unprofitable at every cost level, including at 5~bps, confirming that
the negative performance is not primarily driven by transaction costs. The impact of
higher costs is modest: moving from 5~bps to 25~bps reduces the annualized return by
approximately $140$ basis points and deepens the maximum drawdown by $400$ basis points,
while the Sharpe ratio worsens from $-0.22$ to $-0.28$.

\begin{table}[!ht]
\centering
\caption{Transaction-cost sensitivity. Average monthly one-way turnover $\approx 53.6\%$
         (median $57.2\%$; 95th percentile $83.8\%$). All figures annualized except
         turnover.}
\label{tab:tc}
\begin{tabular}{rrrrrr}
\toprule
Cost (bps) & Ann\,Ret & Ann\,Vol & Sharpe & Max\,DD & Avg\,TO \\
\midrule
\phantom{0}5  & $-$2.44\% & 20.66\% & $-$0.22 & $-$80.2\% & 53.6\% \\
10            & $-$2.79\% & 20.66\% & $-$0.23 & $-$81.2\% & 53.6\% \\
15            & $-$3.14\% & 20.66\% & $-$0.25 & $-$82.3\% & 53.6\% \\
25            & $-$3.84\% & 20.66\% & $-$0.28 & $-$84.2\% & 53.6\% \\
\bottomrule
\end{tabular}
\end{table}

The turnover figures (54\% average one-way per month) are substantially higher than
the 27\% reported in earlier versions of this analysis, which used a defective weight
normalization that halved all portfolio weights. The corrected turnover is consistent
with what one would expect from a strategy that replaces approximately $10$--$20\%$ of
its names each month plus additional drift-induced trading. At 10~bps and 54\% monthly
turnover, annual cost drag is $\approx65$~bps, which is non-trivial but insufficient
to explain the strategy's $\approx240$~bps gross underperformance relative to a
risk-free rate.

\subsection{Rolling Sharpe and Time-Series Behavior}

Figure~\ref{fig:rolling} displays the rolling 36-month Sharpe ratio for the strategy,
the U.S.\ market, and the UMD factor. Three observations stand out. First, the strategy
had a positive rolling Sharpe in the pre-crisis period (2006--early 2008), when momentum
was working well in the trending market environment---consistent with the strategy's
equity curve reaching $1.6\times$ in mid-2008. Second, the rolling Sharpe dropped
sharply negative in 2009 and remained negative or near zero for most of the subsequent
decade as the short leg bled. Third, the strategy's rolling Sharpe closely tracks that
of the UMD factor, with near-identical timing of the 2009 deterioration and the 2020
shock, confirming the factor regression result that the strategy is essentially a
leveraged, cost-burdened exposure to UMD.

\begin{figure}[!ht]
  \centering
  \includegraphics[width=\textwidth]{figures/rolling_sharpe_36m.png}
  \caption{Rolling 36-month Sharpe ratio (minimum 12 months of data) computed as
           annualized mean excess return over the risk-free rate divided by annualized
           standard deviation. Strategy (net) in red; US Market in blue dashed;
           UMD factor in green dot-dash. The close co-movement of the strategy and
           UMD lines confirms that the strategy's performance is primarily UMD exposure.}
  \label{fig:rolling}
\end{figure}

\subsection{Return Distribution}

Figure~\ref{fig:dist} shows the distribution of monthly net strategy returns.
The distribution has a left tail that is substantially fatter than a normal distribution
(skewness $-1.21$, excess kurtosis $5.84$): most months produce modest positive or
near-zero returns, but a small number of months produce catastrophic losses. The two
worst months ($-31\%$ in April~2009, $-19\%$ in November~2020) are visible as isolated
observations far to the left of the bulk of the distribution. This negative skew is a
defining characteristic of momentum strategies noted by \citet{Daniel2016} and is
precisely the feature that makes momentum crashes so damaging: the strategy earns small
positive returns in most periods but periodically suffers losses that wipe out years of
gains.

\begin{figure}[!ht]
  \centering
  \includegraphics[width=\textwidth]{figures/return_distribution.png}
  \caption{Histogram of monthly net strategy returns (226 observations, 50 bins).
           Dashed vertical line indicates the mean return ($-0.26\%$/month).
           Skewness $= -1.21$; excess kurtosis $= 5.84$. The fat left tail reflects
           momentum crash months (April~2009, November~2020).}
  \label{fig:dist}
\end{figure}

\subsection{Factor Regressions}

Table~\ref{tab:ff} reports the results of regressing net excess strategy returns on
four factor models. We discuss the progression of specifications in turn.

Under CAPM, the alpha is $-0.72\%$ annualized with a $t$-statistic of $-0.32$---negative
but statistically indistinguishable from zero. The $R^2$ of $0.147$ indicates that
the market factor explains only $15\%$ of the strategy's return variation, consistent
with the strategy being approximately dollar-neutral (the long and short legs roughly
offset market exposure in normal times).

Adding the size and value factors (FF3) and then the profitability and investment factors
(FF5) modestly increases $R^2$ to $0.27$--$0.28$ but leaves the alpha statistically
insignificant. The strategy has mild positive loadings on SMB (buying slightly smaller
large-caps and shorting slightly larger ones) and negative loading on HML (buying
growth stocks that have had strong momentum and shorting value stocks that have been
beaten down), but these exposures explain limited return variation.

The decisive result comes with the addition of UMD. Including the momentum factor
raises $R^2$ sharply from $0.279$ to $0.823$, indicating that $82\%$ of the strategy's
monthly return variation is spanned by the UMD factor alone. The momentum factor
loading is large and highly significant (not shown in Table~\ref{tab:ff}), confirming
that the strategy is, at its core, a leveraged and cost-burdened implementation of
the UMD factor. The alpha under FF5$+$UMD is $-4.03\%$ annually with $t = -3.23$
(significant at the $1\%$ level), implying that the strategy not only fails to generate
alpha but actually \emph{underperforms} the UMD factor on a risk-adjusted basis once
all five Fama--French factors are controlled for. This underperformance reflects the
friction of transaction costs, the mechanical drag of the $\pm40\%$ return cap relative
to the uncapped UMD factor, and potential implementation differences in universe
definition.

\begin{table}[!ht]
\centering
\caption{Factor regressions, net excess strategy returns. OLS with Newey--West
         HAC standard errors (6 lags). Alpha annualized as $(1 + \hat\alpha_m)^{12} - 1$.
         Sample: March~2006--December~2024 ($n = 226$ months).
         $^{*}p{<}0.10$, $^{**}p{<}0.05$, $^{***}p{<}0.01$.}
\label{tab:ff}
\begin{tabular}{lrrrr}
\toprule
Model & $\hat\alpha$ (ann) & $t(\hat\alpha)$ & $R^2$ & $n$ \\
\midrule
CAPM        & $-$0.72\%          & $-$0.32        & 0.147 & 226 \\
FF3         & $-$1.65\%          & $-$0.70        & 0.269 & 226 \\
FF5         & $-$1.56\%          & $-$0.65        & 0.279 & 226 \\
FF5\,+\,UMD & $-$4.03\%$^{***}$  & $-$3.23        & 0.823 & 226 \\
\bottomrule
\end{tabular}
\end{table}

% ================================================================== %
\section{Discussion}
\label{sec:discussion}
% ================================================================== %

Our results paint a coherent picture that reconciles the apparent contradiction between
a working signal (the long leg earns $+7.9\%$) and a failing strategy ($-2.79\%$ net).
The resolution lies in the asymmetric behavior of the short leg during market dislocations.

In a secular bull market punctuated by sharp reversals, the stocks most likely to be in
the bottom decile of momentum are those that have experienced severe drawdowns---financially
distressed companies, cyclicals caught in a downturn, or sectors experiencing structural
disruption. These stocks carry high beta, high volatility, and high short interest from
the entire market, not just momentum traders. When sentiment shifts (as it did in March--April
2009 and November~2020), these stocks do not merely recover modestly---they snap back
violently, often doubling within weeks. A leveraged short position in such stocks is
effectively writing a far-out-of-the-money put option on the market: it collects a
small premium in most months but occasionally suffers a massive loss when the market
reverses sharply.

This interpretation is consistent with \citet{Daniel2016}'s option-theoretic model of
momentum crashes. They show that momentum strategies have a return profile similar to
a short volatility position: the strategy earns positive carry in trending markets and
suffers acute losses during sharp reversals. Our data confirms this: the rolling Sharpe
(Figure~\ref{fig:rolling}) is positive when the market is trending (2006--2008,
2014--2018) and sharply negative following abrupt reversals.

The implication for implementation is clear. A \emph{long-only} momentum overlay---buying
top-decile momentum stocks as an overweight against an index benchmark---would have earned
approximately $+7.9\%$ gross without the crash exposure. Alternatively,
\citet{Barroso2015}'s volatility-scaling approach---reducing position size when recent
strategy volatility is elevated---has been shown to substantially improve the
risk-adjusted performance of momentum strategies by cutting exposure precisely when
crash risk is highest. Both modifications address the short-leg problem identified here
and represent natural extensions of this research.

% ================================================================== %
\section{Conclusion}
\label{sec:conclusion}
% ================================================================== %

We have examined whether the canonical 12--1 month momentum strategy, applied to the
S\&P~500 with survivorship-bias-free data and realistic transaction costs, generates
abnormal returns over the 2006--2024 period. The answer is clearly negative: the
strategy produces a net annualized return of $-2.79\%$, a Sharpe ratio of $-0.23$, and
a maximum drawdown of $-81.2\%$.

However, our leg decomposition reveals that this negative result masks an important
underlying signal. The long leg---past winners---earns $+7.9\%$ annualized and the
short leg---past losers---loses $-9.1\%$, confirming that momentum ranks stocks
correctly as a cross-sectional predictor. The strategy fails not because the signal
is wrong but because the short leg bears two catastrophic losses---April~2009
($-31\%$) and November~2020 ($-19\%$)---driven by violent snap-back rallies in deeply
distressed stocks. Transaction costs, while non-trivial at $\approx65$~bps annually,
are a secondary concern. Factor regressions confirm that $82\%$ of the strategy's
return variation is explained by the Fama--French UMD factor, with a statistically
significant negative alpha once momentum exposure is controlled for.

These findings have three practical implications. First, the profitability of momentum
in the S\&P~500 depends critically on whether the short leg is implemented and how
crash risk is managed; a long-only implementation of the signal would have been
substantially more profitable over this period. Second, the 2006--2024 window is
an unusually challenging period for large-cap momentum---it contains two of the
largest momentum crashes on record and a prolonged bull market hostile to short-selling
past losers---and results over a longer sample or in smaller-cap universes would likely
differ. Third, the signal's validity as a ranking mechanism, evident in the long leg's
$+7.9\%$ return, suggests that momentum information \emph{is} priced into S\&P~500
stocks, just not in a way that is easily harvested through a symmetric long-short
implementation.

\paragraph{Limitations.}
\begin{enumerate}[label=(\roman*), leftmargin=*]
  \item Membership history relies on Wikipedia; CRSP's historical constituent files
        would provide more complete and accurate coverage, particularly pre-2005.
  \item The $\pm40\%$ monthly return cap is a modeling choice that suppresses measured
        crash severity; without it, April~2009 would register as a $-53\%$ month,
        which more accurately reflects the magnitude of the momentum crash but
        likely exceeds what a real trading desk could have sustained.
  \item Turnover costs are assumed uniform at 10~bps; in practice, liquidity varies
        across names and time periods, and implementation shortfall would likely
        be higher for the most illiquid short positions.
  \item The analysis does not account for stock-borrow costs on the short leg, which
        can be substantial for the most heavily shorted names and would further
        reduce net returns.
\end{enumerate}

\paragraph{Future directions.}
The most natural extension is to implement \citeauthor{Barroso2015}'s (\citeyear{Barroso2015})
volatility-scaling overlay, which has been shown to improve momentum Sharpe ratios by
approximately 50\% by reducing exposure during high-volatility regimes. A second direction
is to expand the universe to include mid- and small-cap stocks, where the momentum
premium has historically been more robust and crash episodes less severe relative to
the signal's overall performance. A third direction is to explore whether the signal's
positive long-leg return ($+7.9\%$) can be harvested cost-effectively through a
long-only factor tilt, potentially avoiding the short-leg crash risk entirely.

\paragraph{Reproducibility.}
All code, data pipelines, and scripts to replicate every table and figure in this paper
are publicly available at \url{https://github.com/dshan12/Momentum-Research}. The
analysis is fully automated: running the pipeline from raw Wikipedia and Yahoo Finance
data through to the final PDF requires no manual intervention.

% ================================================================== %
\bibliographystyle{apalike}
\bibliography{references}

\end{document}
